\section{结语：江海能力并不取消边界的代价}

南宋的持续存在，说明东南的财政、交通、市场与地方组织能够在失去北方后支撑一个国家；但这种能力不是无边界的富裕。和议需要军费和转运，海贸依赖港口管理与航路风险，纸币和税粮又把中枢的需求转化为各地不同的折算与负担。金、南宋与蒙古的行动持续改变这些安排，故守势、北伐与海域经营都应放回多政权的资源竞争中理解。

一二七九年结束的是南宋朝廷，不是江南的社会与制度。城市、港口、士人、军户和地方家庭以不等速的方式进入元的接收秩序，也保留了先前形成的交通、知识与经济关系。蒙古征服一章将讨论这种接收怎样横跨西夏、金与南宋旧地，而非把它写成单一军队越过地图的故事。
